\begin{figure}[htbp]
\centering
\begin{tikzpicture}[font=\small,>=Latex,line cap=round,line join=round]
  \begin{scope}[shift={(-5.3,0)}]
    \coordinate (fa) at (0,0);\coordinate (fb) at (3.2,0);
    \coordinate (fc) at (2.7,2);\coordinate (fd) at (.3,1.5);
    \draw[mpblue!15] (1.6,.6625) circle (1.731735);
    \fill[mpblue!5] (fa)--(fb)--(fc)--cycle;
    \fill[mpteal!5] (fa)--(fc)--(fd)--cycle;
    \draw[mpgray] (fa)--(fb)--(fc)--(fd)--cycle;
    \draw[mpochre,line width=1.35pt] (fa)--(fc);
    \foreach \p in {fa,fb,fc,fd}\fill[mpblue] (\p) circle (1.8pt);
    \node[below left] at (fa) {$A$};\node[below right] at (fb) {$B$};
    \node[above right] at (fc) {$C$};\node[above left] at (fd) {$D$};
    \node[font=\bfseries\small] at (1.6,2.85) {The diagonal $AC$};
    \node[text=mpochre,align=center,font=\scriptsize] at (1.6,-1.35)
      {Not locally Delaunay\\$D$ lies inside the circle through $A,B,C$};
  \end{scope}
  \draw[->,mpgray,line width=.85pt] (-1.1,.85)--(.9,.85)
    node[midway,above] {flip};
  \begin{scope}[shift={(1.8,0)}]
    \coordinate (ga) at (0,0);\coordinate (gb) at (3.2,0);
    \coordinate (gc) at (2.7,2);\coordinate (gd) at (.3,1.5);
    \fill[mpblue!5] (ga)--(gb)--(gd)--cycle;
    \fill[mpteal!5] (gb)--(gc)--(gd)--cycle;
    \draw[mpgray] (ga)--(gb)--(gc)--(gd)--cycle;
    \draw[mpteal,line width=1.35pt] (gb)--(gd);
    \foreach \p in {ga,gb,gc,gd}\fill[mpblue] (\p) circle (1.8pt);
    \node[below left] at (ga) {$A$};\node[below right] at (gb) {$B$};
    \node[above right] at (gc) {$C$};\node[above left] at (gd) {$D$};
    \node[font=\bfseries\small] at (1.6,2.85) {The diagonal $BD$};
    \node[text=mpteal,align=center,font=\scriptsize] at (1.6,-1.35)
      {Locally Delaunay\\the planar vertices have not moved};
  \end{scope}
  \node[align=center] at (0,-2.35)
    {For the same nodal values $u$: \quad
     $D_{\mathcal T_{BD}}(u)\le D_{\mathcal T_{AC}}(u)$.};
\end{tikzpicture}
\caption{A diagonal flip on a fixed Euclidean quadrilateral, shown
with its actual vertex positions. Rippa's energy comparison selects the
Delaunay diagonal for fixed nodal values. The same combinatorial move
also occurs in flip graphs used in surface topology; the displayed
energy comparison concerns the fixed Euclidean realization.}
\label{fig:mono-flip}
\end{figure}
